Substitute $H_i = p_i/\tau + \sum_{n\geq 1}h_{i,n}\tau^{2n-1}$. At order $\tau^{-2}$ the evolution equations give $-p_i + p_i\sum_j p_j = 0$ and the constraint gives $\sum_{i<j}p_ip_j = 0$: the Kasner relations. At order $\tau^{2n-2}$, $n \geq 1$, the evolution equation for $i$ reads $(2n-1)h_{i,n} + h_{i,n}\sum_j p_j + p_i\sum_j h_{j,n} = \Lambda\,\delta_{n,1} - K_{i,n}$, where $K_{i,n}$ collects the products of coefficients already determined; with $\sum p = 1$ this is $2n\,h_{i,n} + p_i\,\mathcal{H}_n = \Lambda\delta_{n,1} - K_{i,n}$ with $\mathcal{H}_n := \sum_j h_{j,n}$. Summing over $i$ gives $(2n+1)\mathcal{H}_n = 3\Lambda\delta_{n,1} - \sum_iK_{i,n}$, which determines $\mathcal{H}_n$ and then each $h_{i,n}$ uniquely. For $n = 1$, $K_{i,1} = 0$, so $\mathcal{H}_1 = \Lambda$ and $h_{i,1} = (\Lambda - p_i\Lambda)/2$. Odd powers only: no even power of $\tau$ is sourced, since the equations are quadratic in $H$ with the pole at odd order and $\Lambda$ at order zero. The constraint is preserved by the evolution: with $S := \sum_iH_i$ and $C := \sum_{i<j}H_iH_j$ the evolution equations give $\dot C = 2S(\Lambda - C)$, so $C - \Lambda$ obeys a linear homogeneous equation whose solutions near $\tau = 0$ behave as $\tau^{-2}$, and the Kasner relation $\sum_{i<j}p_ip_j = 0$ sets that mode's coefficient to zero; the constraint therefore holds to all orders once it holds at the pole. Odd powers only is likewise forced: an even power $c_i\tau^{2n}$ in $H_i$, $n \geq 0$, would enter the evolution equations at order $\tau^{2n-1}$, where nothing sources it, through three terms, $2n\,c_i$ from $\dot H_i$, $c_i\sum_jp_j = c_i$ from $c_i\tau^{2n}$ against the pole of $\sum_jH_j$, and $p_i\sum_jc_j$ from the pole of $H_i$ against $\sum_jc_j\tau^{2n}$; the system $(2n+1)c_i + p_i\sum_jc_j = 0$ sums to $(2n+2)\sum_jc_j = 0$, hence $\sum_jc_j = 0$ and $c_i = 0$; the script confirms this by direct substitution to order $\tau^4$ at the Schwarzschild exponents with $\Lambda = 3$, where $h_{i,1} = (2, \tfrac12, \tfrac12)$, $h_{i,2} = (-\tfrac{33}{20}, -\tfrac{3}{40}, -\tfrac{3}{40})$, and both residuals vanish identically through the order computed. The Schwarzschild-de Sitter interior is not in this class (its slices are Kantowski-Sachs); what the geometry paper's Theorem 5.9 \cite[Thm.~5.9]{ButlerT6} (\claimref{interior.schwarzschild-de-sitter-weighted-limit}) gives is that it shares the residue and the exponents of the Schwarzschild approach, with the same weighted limit and rates (\claimref{dark-energy.imported-weighted-limit-schwarzschild-singularity-matching}), which is all the statement claims for it.
